% ══════════════════════════════════════════════════════════════════════════
%  Chapter 3: Influence, Noise Sensitivity, and the Noise Operator
% ══════════════════════════════════════════════════════════════════════════

\chapter{Influence and Noise Sensitivity}\label{ch:influence-noise}

This chapter introduces three related concepts---influence, noise
sensitivity, and the noise operator---that quantify how ``sensitive'' a
Boolean function is to perturbations of its input. Each has a clean
spectral characterization, linking combinatorial/probabilistic notions to
the Fourier coefficients. The main reference is
\citet[Chapters~2--3]{odonnell2014analysis}.

% ──────────────────────────────────────────────────────────────────────
\section{Influence of Variables}\label{sec:influence}

Informally, the \emph{influence} of variable $i$ measures how often
flipping $x_i$ changes the output of $f$.

\begin{definition}[Influence]\label{def:influence}
For $f \colon \pmone^n \to \pmone$ and $i \in [n]$, the \emph{influence
of variable $i$} is
\begin{equation}\label{eq:influence-def}
  \Inf_i[f] = \Prob_x[f(x) \neq f(x^{\oplus i})],
\end{equation}
where $x^{\oplus i}$ denotes $x$ with the $i$-th coordinate flipped:
$(x_1, \ldots, -x_i, \ldots, x_n)$.
\end{definition}

\begin{theorem}[Spectral formula for influence]\label{thm:influence-spectral}
\begin{equation}\label{eq:influence-spectral}
  \Inf_i[f] = \sum_{S \ni i} \fhat(S)^2.
\end{equation}
\end{theorem}

\begin{proof}
Define $D_i f(x) = \frac{1}{2}(f(x) - f(x^{\oplus i}))$, the
\emph{discrete derivative} in direction $i$. Then
$f(x) \neq f(x^{\oplus i})$ iff $|D_i f(x)| = 1$ (since $f$ takes
values $\pm 1$, the difference is $0$ or $\pm 2$). Thus:
\[
  \Inf_i[f] = \E_x[(D_i f(x))^2].
\]
Now compute the Fourier expansion of $D_i f$. Since
$\chi_S(x^{\oplus i}) = (-1)^{[i \in S]} \chi_S(x)$ (flipping $x_i$
negates $\chi_S$ iff $i \in S$):
\[
  D_i f(x) = \tfrac{1}{2}\sum_S \fhat(S)\,\chi_S(x)\,(1 - (-1)^{[i \in S]})
  = \sum_{S \ni i} \fhat(S)\,\chi_S(x).
\]
By Parseval: $\E[(D_i f)^2] = \sum_{S \ni i} \fhat(S)^2$.
\end{proof}

\begin{definition}[Total influence]\label{def:total-influence}
\begin{equation}\label{eq:total-influence}
  \mathbf{I}[f] = \sum_{i=1}^n \Inf_i[f]
  = \sum_{S \subseteq [n]} |S| \cdot \fhat(S)^2.
\end{equation}
\end{definition}

The total influence is the expected ``degree'' under the spectral
distribution $\{\fhat(S)^2\}_{S \subseteq [n]}$. Functions with
low-degree spectral concentration necessarily have low total influence,
and conversely.

\begin{example}[Influences of example functions]\label{ex:influences}
\begin{itemize}
  \item \textbf{Dictator $f(x) = x_1$:}
    $\Inf_1[f] = 1$, $\Inf_i[f] = 0$ for $i > 1$. Total: $\mathbf{I}[f] = 1$.
  \item \textbf{Parity $f(x) = x_1 \cdots x_n$:}
    $\Inf_i[f] = 1$ for all $i$. Total: $\mathbf{I}[f] = n$.
  \item \textbf{$\text{Maj}_3$:}
    $\Inf_i[f] = \fhat(\{i\})^2 + \fhat(\{1,2,3\})^2 = 1/4 + 1/4 = 1/2$
    for each $i$. Total: $\mathbf{I}[f] = 3/2$.
\end{itemize}
\end{example}

The total influence is bounded: $\mathbf{I}[f] \leq n$ for any Boolean
$f$ (since each $\Inf_i[f] \leq 1$). The maximum is achieved by the
parity function. A celebrated theorem of
\citet{kahn1988influence} shows that every Boolean function depending on
all $n$ variables has $\max_i \Inf_i[f] \geq \Omega(\log n / n)$.


% ──────────────────────────────────────────────────────────────────────
\section{The Noise Operator}\label{sec:noise-operator}

\begin{definition}[Noise operator]\label{def:noise-operator}
For $\rho \in [-1, 1]$, the \emph{Bonami--Beckner noise operator}
$T_\rho$ acts on $f \colon \pmone^n \to \R$ by
\begin{equation}\label{eq:noise-operator}
  T_\rho f(x) = \E_y[f(y)],
\end{equation}
where $y$ is the ``$\rho$-correlated copy'' of $x$: each coordinate
$y_i$ independently equals $x_i$ with probability $\frac{1+\rho}{2}$
and $-x_i$ with probability $\frac{1-\rho}{2}$.
\end{definition}

Equivalently, $y_i = x_i \cdot z_i$ where $z_i$ is an independent
$\pm 1$ random variable with $\E[z_i] = \rho$.

The crucial property of the noise operator is its diagonal action in the
Fourier basis:

\begin{theorem}[Spectral characterization of $T_\rho$]\label{thm:noise-spectral}
\begin{equation}\label{eq:noise-spectral}
  T_\rho f(x) = \sum_{S \subseteq [n]} \rho^{|S|}\, \fhat(S)\, \chi_S(x).
\end{equation}
That is, $T_\rho$ multiplies each Fourier coefficient of degree $k$ by
$\rho^k$.
\end{theorem}

\begin{proof}
It suffices to check on the basis elements $\chi_S$:
\begin{align*}
  T_\rho \chi_S(x)
  &= \E_y\!\left[\prod_{i \in S} y_i\right]
  = \E_y\!\left[\prod_{i \in S} x_i z_i\right]
  = \chi_S(x) \cdot \prod_{i \in S} \E[z_i]
  = \rho^{|S|}\, \chi_S(x).
\end{align*}
(The $z_i$ are independent, so the expectation factors.)
\end{proof}

\begin{remark}[Interpretation]\label{rem:noise-interpretation}
For $\rho$ close to~1, $T_\rho$ barely changes $f$: it slightly damps
high-degree coefficients. For $\rho = 0$, $T_\rho f(x) = \fhat(\emptyset)$
is just the mean. For $\rho$ close to~0, only the low-degree part of $f$
survives. The noise operator is a ``low-pass filter'' on the Boolean
hypercube.
\end{remark}


% ──────────────────────────────────────────────────────────────────────
\section{Noise Stability}\label{sec:noise-stability}

\begin{definition}[Noise stability]\label{def:noise-stability}
The \emph{noise stability} of $f$ at correlation $\rho$ is
\begin{equation}\label{eq:stab-def}
  \Stab_\rho[f] = \ip{f}{T_\rho f}
  = \E_{x,y}[f(x)\, f(y)],
\end{equation}
where $(x,y)$ is a $\rho$-correlated pair.
\end{definition}

\begin{theorem}[Spectral formula]\label{thm:stab-spectral}
\begin{equation}\label{eq:stab-spectral}
  \Stab_\rho[f] = \sum_{S \subseteq [n]} \rho^{|S|}\, \fhat(S)^2.
\end{equation}
\end{theorem}

\begin{proof}
By Plancherel (\Cref{thm:plancherel}):
$\ip{f}{T_\rho f} = \sum_S \fhat(S) \cdot \rho^{|S|} \fhat(S) = \sum_S \rho^{|S|} \fhat(S)^2$.
\end{proof}

\begin{definition}[Noise sensitivity]\label{def:noise-sensitivity}
The \emph{noise sensitivity} is
\[
  \mathbf{NS}_\delta[f] = \Prob_{(x,y)}[f(x) \neq f(y)],
\]
where $y$ is obtained from $x$ by flipping each bit independently with
probability $\delta$, i.e., $\rho = 1 - 2\delta$. For Boolean $f$:
\[
  \mathbf{NS}_\delta[f] = \frac{1 - \Stab_{1-2\delta}[f]}{2}.
\]
\end{definition}

\paragraph{Relevance to the thesis.}
Noise stability quantifies robustness: a layer with high noise stability
tolerates occasional bit flips in its weights. The thesis uses noise
stability as a spectral diagnostic and a regularization target (the
``stability reward'' in Section~4.5 of the thesis encourages low-degree
concentration by maximizing $\Stab_\rho[f]$).


% ──────────────────────────────────────────────────────────────────────
\section{Relationship Between Influence and Noise Sensitivity}\label{sec:inf-ns}

Total influence and noise sensitivity are connected:

\begin{proposition}\label{prop:inf-ns}
$\mathbf{I}[f] = \lim_{\delta \to 0} \frac{\mathbf{NS}_\delta[f]}{\delta}$.
\end{proposition}

\begin{proof}
$\mathbf{NS}_\delta[f] = \frac{1}{2}(1 - \sum_S (1-2\delta)^{|S|} \fhat(S)^2)$.
Differentiating at $\delta = 0$:
$\frac{d}{d\delta}\mathbf{NS}_\delta[f]\big|_{\delta=0} = \sum_S |S| \cdot \fhat(S)^2 = \mathbf{I}[f]$.
\end{proof}

Total influence is the infinitesimal rate of change of noise sensitivity.
Functions with high total influence are very sensitive to small
perturbations; functions with low total influence are noise-robust.

\begin{example}[Noise stability of example functions]\label{ex:noise-stab}
\begin{itemize}
  \item \textbf{Dictator:} $\Stab_\rho[x_1] = \rho$.
  \item \textbf{Parity:} $\Stab_\rho[x_1 \cdots x_n] = \rho^n \to 0$
    exponentially. Parity is extremely noise-sensitive.
  \item \textbf{$\text{Maj}_3$:} $\Stab_\rho = 3 \cdot \frac{1}{4}\rho + \frac{1}{4}\rho^3 = \frac{3\rho + \rho^3}{4}$.
\end{itemize}
\end{example}
